\documentclass[11pt,a4paper]{article}

\usepackage[utf8]{inputenc}
\usepackage[T1]{fontenc}
\usepackage{lmodern}
\usepackage{amsmath,amssymb,amsthm,mathtools}
\usepackage{geometry}
\geometry{margin=2.7cm}
\usepackage{hyperref}
\usepackage{enumitem}
\usepackage{xcolor}
\usepackage{booktabs}
\usepackage{tcolorbox}
\tcbuselibrary{breakable,skins}
\usepackage{microtype}

\newtheorem{theorem}{Theorem}[section]
\newtheorem{lemma}[theorem]{Lemma}
\newtheorem{proposition}[theorem]{Proposition}
\newtheorem{corollary}[theorem]{Corollary}
\newtheorem{conjecture}[theorem]{Conjecture}
\theoremstyle{definition}
\newtheorem{definition}[theorem]{Definition}
\newtheorem{remark}[theorem]{Remark}
\newtheorem{principle}[theorem]{Diagnostic thesis}

\newcommand{\Q}{\mathbb{Q}}
\newcommand{\Z}{\mathbb{Z}}
\newcommand{\R}{\mathbb{R}}
\newcommand{\C}{\mathbb{C}}
\newcommand{\A}{\mathbb{A}}
\newcommand{\F}{\mathbb{F}}

% status tags (shared with the capstone)
\newcommand{\TAGthm}{{\small\textsf{[\textcolor{green!45!black}{theorem}]}}}
\newcommand{\TAGver}{{\small\textsf{[\textcolor{blue!60!black}{verified}]}}}
\newcommand{\TAGev}{{\small\textsf{[\textcolor{orange!80!black}{evidence}]}}}
\newcommand{\TAGconj}{{\small\textsf{[\textcolor{purple}{conjecture}]}}}
\newcommand{\TAGopen}{{\small\textsf{[\textcolor{red!70!black}{open}]}}}
\newcommand{\TAGphil}{{\small\textsf{[\textcolor{gray}{interpretation}]}}}

\hypersetup{colorlinks=true, linkcolor=blue!55!black,
  citecolor=blue!55!black, urlcolor=blue!55!black}

\title{\textbf{The Void and the Structure}\\
\large What the Riemann zeros are, how the primes grow from them,\\
and why the problem resists construction --- a diagnosis, not an impossibility}
\author{Lee Smart \\ \small Institute of Vibrational Field Dynamics}
\date{2026-05-30 \quad v0.1 (draft)}

\begin{document}
\maketitle

\begin{center}
\fbox{\parbox{0.93\textwidth}{\centering\small
\textsc{Pre-peer-review open research essay.} \textbf{This paper proves
neither the Riemann Hypothesis nor its unprovability.} It does two honest
things: (i) it describes the nontrivial zeros as a generative \emph{absence}
and states precisely, with classical results, how the primes are
reconstructed from them; and (ii) it offers a \emph{diagnosis of difficulty}
--- a principled account of why RH has resisted every constructive method.
Claims are tagged \TAGthm, \TAGver, \TAGev, \TAGconj, \TAGopen, or
\TAGphil\ (a marked, non--load-bearing interpretive register). The
diagnosis is a heuristic thesis (\TAGconj), \emph{not} a theorem; in
particular we do \textbf{not} claim RH can never be proved (see
\S\ref{sec:boundary}).}}
\end{center}

\begin{abstract}
The nontrivial zeros of an L-function are not a constructed object but an
\emph{absence}: the locus where the function vanishes, realised in Connes'
spectral picture as an \emph{absorption spectrum} --- missing lines in a
continuous scaling spectrum. We take this seriously as the object's defining
character. First, we record the classical fact that the entire multiplicative
structure of the integers --- the primes --- is reconstructed \emph{from} this
absence: the Riemann--von Mangoldt explicit formula writes the prime-counting
staircase as a sum of waves whose frequencies are exactly the zeros. The
structure grows out of the void. Second, we describe the void's
character as far as proof allows: it is eternal in the precise sense of being
timeless, uncreated, and observer-independent; it is unbounded; and beyond a
fixed set of arithmetic constraints it is \emph{formless} in an exact sense
(its fine statistics are those of the most structureless Hermitian ensemble,
GUE). Third --- the technical heart --- we give a diagnosis of why RH resists
construction: random-matrix \emph{universality} hands you the void's
\emph{class} for free, but its specific member is fixed only by the primes, so
any construction that does not already encode the arithmetic underdetermines
the object to the universal class. Proof proceeds by building structure; the
target is an absence pinned only by the arithmetic. That mismatch, we argue,
is the friction. We are explicit that this is a diagnosis of \emph{hardness},
not a proof of \emph{impossibility}: no independence result is known, much is
known about the void, and were RH ever shown independent of a sound theory it
would thereby be true.
\end{abstract}

\tableofcontents

\section{Stance}\label{sec:stance}
This is a companion to the frontier capstone of this bundle. Its computations
are those of the capstone, together with one demonstration new to this
companion (Prop.\ \ref{prop:detprob}, \texttt{deterministic\_probability.py});
no result here exceeds the capstone's verified status. Its purpose is
conceptual: to say cleanly \emph{what} the zeros are, \emph{how} the primes
issue from them, and \emph{why} the problem is hard. Two disciplines hold
throughout. (1) The mathematics stays mathematical: every load-bearing
statement is a theorem, a verified computation, or an explicitly named open
problem. (2) Where the essay reaches for words like ``eternal'' or
``timeless,'' those words are used in a precise, defensible sense and are
marked \TAGphil; any further meaning a reader brings is theirs, and the text
neither asserts nor needs it (\S\ref{sec:discipline}).

\section{The void: an absence, not an object}\label{sec:void}

\begin{definition}[The void]\label{def:void}
For an L-function $L(s)$, the \emph{void} is the set
$Z=\{\rho : L(\rho)=0,\ 0<\Re\rho<1\}$ of nontrivial zeros --- the points where
$L$ \emph{takes the value zero}. (The concrete claims below --- primes,
$\psi(x)$, $N(T)$, the proportion on the line, RH --- are stated for the
Riemann zeta function; the ``$L$-function'' phrasing marks the general frame,
and we flag where generality is not yet established.)
\end{definition}

The word ``absence'' below refers not to the analytic value at $\rho$ (which is
$0$), but to the \emph{spectral} reading: in Connes' framework the zeros are
the lines \emph{missing} from a continuous spectrum. Two features make
``absence,'' in that sense, the right word rather than a metaphor.

\begin{proposition}[Spectral realisation as absorption, \TAGthm\ (Connes
\cite{Connes})]\label{prop:absorption}
On the adele class space $\A_\Q/\Q^\times$ the idele-class group acts by
scaling; the zeros appear not as eigenvalues added to a spectrum but as an
\emph{absorption spectrum} --- the frequencies \emph{missing} from an otherwise
continuous spectrum. RH is the statement that these absences lie on the
critical line.
\end{proposition}

\begin{remark}[Why ``not an object'']\label{rem:notobject}
The Hilbert--Pólya programme seeks a self-adjoint operator whose spectrum
\emph{is} $Z$. No such operator has been constructed (\TAGopen). Until it is,
$Z$ is specified by a \emph{condition} ($L=0$), not by a construction; and in
the one rigorous spectral picture we have (Proposition
\ref{prop:absorption}), $Z$ is precisely what is \emph{absent}. The void is
something you \emph{locate}, not something you \emph{build}. This is not a
deficiency of effort; \S\ref{sec:resist} argues it is intrinsic to the object.
\end{remark}

\section{The structure grows from the void}\label{sec:grows}

The decisive fact --- and the honest form of ``the void generates the
structure'' --- is classical.

\begin{theorem}[Explicit formula; Riemann, von Mangoldt, \TAGthm]\label{thm:explicit}
For $x>1$ not a prime power, with $\psi(x)=\sum_{p^k\le x}\log p$ the
Chebyshev prime-counting function,
\[
  \psi(x) \;=\; x \;-\; \sum_{\rho}\frac{x^{\rho}}{\rho}
                  \;-\; \log(2\pi) \;-\; \tfrac12\log\!\bigl(1-x^{-2}\bigr),
\]
the sum over nontrivial zeros $\rho$ taken in the usual symmetric limiting
sense ($\lim_{T\to\infty}\sum_{|\Im\rho|\le T}$). Writing
$\rho=\tfrac12+i\gamma$ (the on-line case), the zero term is a sum of
\emph{waves},
\[
  \sum_{\gamma>0}\frac{2\,x^{1/2}}{|\rho|}\,
     \cos\!\bigl(\gamma\log x-\arg\rho\bigr),
\]
each zero contributing one frequency $\gamma$ in the variable $\log x$.
\end{theorem}

\begin{corollary}[\TAGthm]\label{cor:fromvoid}
For $\zeta$, the prime-counting function $\psi(x)$ is reconstructed from the
nontrivial zeros together with the standard smooth and trivial-zero terms: it
is $x$ corrected by the harmonics of the zeros. Since prime powers, hence the
primes, are recoverable from $\psi$, the prime distribution is in this sense
the \emph{interference pattern} of the void's frequencies.
\end{corollary}

\begin{remark}[Diffraction; the two sides of one identity, \TAGver]
The same identity, read the other way, is a diffraction: the formal spectrum
$D(t)=\sum_\gamma\cos(\gamma t)$ has its sharp peaks (Bragg peaks) exactly at
$t=\log p^k$. The identity is classical (Dyson's quasicrystal \cite{Dyson});
the script \texttt{route\_b/diffraction\_quasicrystal.py} illustrates it
($11/11$ peaks matched), it does not establish it. Zeros diffract into primes; primes diffract into zeros. Neither
is prior --- but if one insists on a direction of generation, Theorem
\ref{thm:explicit} writes the \emph{primes as built from the zeros}.
\end{remark}

So ``the geometry grows out of the void'' has an exact reading, and it stops
exactly where the mathematics stops: the void generates the \emph{primes}.
Any extension of ``generates'' beyond arithmetic is interpretation
(\S\ref{sec:discipline}), not a consequence of Theorem \ref{thm:explicit}.

\section{The character of the void}\label{sec:character}

We can say a great deal about the void --- which is itself the refutation of
``nothing can be known about it'' (\S\ref{sec:boundary}). What follows is
defensible in the marked sense.

\begin{proposition}[It is populated and counted, \TAGthm]\label{prop:counted}
There are infinitely many zeros (Hardy); the count to height $T$ is
$N(T)=\frac{T}{2\pi}\log\frac{T}{2\pi}-\frac{T}{2\pi}+O(\log T)$
(Riemann--von Mangoldt). A positive proportion --- at least $2/5$
\cite{Conrey} --- lie on the line, and the lowest zeros have been verified on
the line to very large height \cite{Platt,Gourdon}. The void is partially but
extensively charted; what is open is a single uniform statement about
\emph{all} of it.
\end{proposition}

\begin{proposition}[Zero correlations match GUE under restricted support,
\TAGthm/\TAGev]\label{prop:formless}
For test functions of suitably restricted support, the pair correlation
(Montgomery \cite{Montgomery}, under the stated conventions) and the $n$-level
correlations (Rudnick--Sarnak \cite{RudnickSarnak}) of the zeros agree with
those of the \textbf{GUE} ensemble. This is a theorem \emph{within those
support restrictions and hypotheses}; the unrestricted GUE law for all
statistics is a conjecture. Numerically, every test performed in this bundle
is consistent with GUE (\texttt{route\_b/quantum\_chaos\_test.py},
\texttt{shape\_search.py}, \texttt{deterministic\_probability.py}), with no
counterexample found.
\end{proposition}

\begin{remark}[``Formless,'' read honestly, \TAGphil]\label{rem:formless}
GUE is the relevant \emph{unitary} random-matrix universality class. Calling
the void ``formless'' is an interpretive gloss on Prop.\ \ref{prop:formless}:
to the extent the statistics have been settled they show \emph{no structure
beyond} the unitary universality class, and no additional structure has been
found where they have not. The word is a useful summary, not a theorem that
the void has ``no structure beyond arithmetic constraints'' --- that stronger
statement is not proved.
\end{remark}

\begin{proposition}[It is timeless, uncreated, observer-free, \TAGphil]\label{prop:eternal}
$Z$ is a property of $L$, hence of the integers: it is not located in space or
in time, it was not brought into being and cannot be unmade, and it is
identical for every reasoner who computes it. In this precise sense --- and
\emph{only} this sense --- the word ``eternal'' is exact. It requires no belief:
the zeros are reached by proof and computation, not faith.
\end{proposition}

\begin{remark}[How a formless object is known, \TAGphil]\label{rem:knowing}
``Formless'' (Prop.\ \ref{prop:formless}) does not mean
``informationless.'' It means the object is known by its \emph{law}, not by
a \emph{figure} --- exactly as a gas has no shape yet an exact
thermodynamics (temperature, pressure, the Maxwell--Boltzmann law) known
precisely and \emph{provably}. The information a formless object yields is
statistical and structural, and it can be exact. For the void, moreover, much
of this knowledge is \emph{proof}, not inference: the count $N(T)$, the
proportion of zeros on the line, the explicit formula (Theorem
\ref{thm:explicit}), and the restricted-support correlations (Prop.\
\ref{prop:formless}, within their support restrictions) are \emph{theorems},
deduced, not measured. What remains \emph{inference} --- the step from GUE
statistics to a self-adjoint generator (the Hilbert--Pólya picture) --- is held
separately as \emph{evidence} and never as proof, and that is precisely the
seam at which RH is still open. So the honest position answers the natural
worry directly: a formless object is knowable \emph{through its law} --- partly
by rigorous theorems, partly by tested statistical prediction; the parts we
have proved are undeniable; and the single part that is only \emph{inferred} is
exactly the part that keeps RH open. The physical analogue is real, not rhetorical: GUE statistics are
\emph{measured} in heavy nuclei and chaotic microwave cavities (Wigner;
Montgomery--Odlyzko) --- one observes a formless ensemble through its level
spacings, never by viewing a shape. Formlessness is a limit on \emph{figure},
not on \emph{knowledge}.
\end{remark}

\section{The void's law is the quantum-chaos law}\label{sec:quantum}

The zeros' statistics (Prop.\ \ref{prop:formless}) are those expected for
\textbf{GUE}, the law of quantum chaotic spectra with broken time-reversal
symmetry --- proven under restricted support and matched numerically; the same
law measured in heavy nuclei \cite{Wigner} and chaotic microwave cavities
\cite{Stockmann}. We make the resulting ``deterministic probability'' explicit.

\begin{proposition}[Deterministic object, quantum-chaos statistics, \TAGev,
\texttt{route\_b/deterministic\_probability.py}]\label{prop:detprob}
The first $300$ Riemann zeros, computed numerically to high precision
(\texttt{mpmath.zetazero}; a fixed, deterministic arithmetic object), have
consecutive-spacing ratio $\langle r\rangle=0.617$; a \emph{genuinely random}
GUE matrix ($600{\times}600$, bulk eigenvalues, fixed seed) gives $\langle
r\rangle=0.593$. Both lie near the GUE value $0.6027$ and far from Poisson
($0.386$) and from a rigid lattice ($1.0$). The statistic $\langle r\rangle$
is unfolding-free, so the comparison uses \emph{no} fitted parameter. (One
matrix draw and $300$ zeros are a finite sample, not a stable benchmark; the
point is qualitative agreement of class, not a precise constant.)
\end{proposition}

\begin{remark}[Necessity and chance in one object, \TAGphil]
The zeros are not random --- $\zeta$ is one fixed function. Yet they exhibit
the statistics predicted for a quantum-chaotic (GUE) spectrum. This is
\emph{deterministic pseudorandomness} (Montgomery--Odlyzko \cite{Odlyzko};
established, reproduced here, not new): the void is where \emph{necessity}
(the zeros are determined) and \emph{chance} (they wear the GUE law) coincide
in the same object. This is evidence for the spectral (Hilbert--Pólya)
picture; it does \emph{not} exhibit the operator and does \emph{not} advance
RH.
\end{remark}

\begin{remark}[Consilience with the programme's Born-rule work, \TAGphil]
The same principle --- a deterministic substrate that appears probabilistic ---
is proposed physically in this programme's crystallisation papers
\cite{BridgePaper,ProbGeom}, where a deterministic constraint-selection
mechanism is conjectured to underlie the Born rule $|\langle\psi_i|\psi
\rangle|^2$. The bridge paper \cite{BridgePaper} states plainly that its
Born-rule recovery is \emph{conjectured, not proven}; Paper~XXX
\cite{ProbGeom} argues a \emph{conditional} reconstruction under imported
assumptions --- a weaker-than-unconditional but stronger-than-conjecture
status. The connection to the void is a \emph{shared principle}, not a
derivation, and across \emph{different} notions of probability: the Born rule
concerns measurement \emph{outcomes}, whereas the void concerns \emph{spectral
statistics}. The void's modest contribution is a clean, parameter-free
\emph{illustration} that the principle can hold --- a deterministic object
whose statistics are, in the tested and restricted senses, consistent with a
quantum-random one. It does not derive the Born rule and does not discharge
the open conjecture of \cite{BridgePaper}.
\end{remark}

\section{Why the problem resists construction}\label{sec:resist}

We now give the diagnosis. It is a heuristic thesis (\TAGconj), supported by
the structure established above and by the explicit computations of this
bundle, \emph{not} a theorem.

\begin{proposition}[Universality gives the class for free, \TAGev]\label{prop:universal}
GUE statistics are \emph{universal}: broad classes of Hermitian random-matrix
ensembles share them under precise universality hypotheses (the general
mathematical fact, independent of the zeros). In this bundle, a rigid scaling
generator perturbed by a \emph{generic} Hermitian term moves to the GUE class
(\texttt{route\_b/first\_principles\_scaling.py}), and a search over shapes
matches the zeros to the GUE \emph{ensemble} effortlessly while no symmetric or
polytopal figure fits (\texttt{shape\_search.py}). The \emph{class} of the void
is cheap to reproduce.
\end{proposition}

\begin{proposition}[The member is fixed by the arithmetic, \TAGthm/\TAGev]
\label{prop:member}
By Theorem \ref{thm:explicit} the specific zeros are in exact duality with the
specific primes. The GUE \emph{ensemble} is shared by countless operators; the
\emph{one} member whose spectrum is $Z$ is singled out, in the
explicit-formula sense, by the full arithmetic data of $\Z$.
\end{proposition}

\begin{principle}[The construction--absence mismatch, \TAGconj]\label{prin:mismatch}
The known routes to RH proceed by \emph{constructing a structure}: a
self-adjoint operator (Hilbert--Pólya, Berry--Keating), a positivity
(Weil/Connes, de Branges, Li), or a geometry (a curve over $\F_1$, a
foliated space). Heuristically, a construction that does not already encode the
full prime data reproduces (by universality, Prop.\ \ref{prop:universal}) the
right \emph{class} but not the right \emph{member} (Prop.\ \ref{prop:member});
to reach the member one must feed in the arithmetic, and
making \emph{that} self-adjoint (or positive) \emph{without already assuming
the conclusion} is the recurrent wall. In short: \textbf{our methods define by
what they build (presence); the target is an absence carrying, as far as is
known, no structure beyond its universality class. The friction between a
constructive method and a formless, arithmetically-pinned target is, on this
diagnosis, why RH is hard.}
\end{principle}

\begin{remark}[``You need shape to define something'']
This is the exact, defensible form of that intuition. A proof needs a
\emph{handle} --- an operator, a form, a variety --- something with shape. The
void offers an absence whose only ``shape'' beyond the arithmetic is maximal
formlessness (Proposition \ref{prop:formless}). The intuition is right as a
\emph{diagnosis of difficulty}; it is \emph{not} a proof that no handle exists
(\S\ref{sec:boundary}). The honest frontier --- the $\F_1$ / arithmetic-site
programme (Connes--Consani; Deninger) --- is precisely the attempt to build a
geometry in which the void acquires a handle and Weil's function-field proof
transfers.
\end{remark}

\section{The boundary: hardness is not impossibility}\label{sec:boundary}

This section is the integrity of the paper. The diagnosis of
\S\ref{sec:resist} explains why RH is \emph{hard}; it does \emph{not} show RH
is \emph{unprovable}, and we explicitly decline that stronger claim, for three
independent reasons.

\begin{enumerate}[label=(\arabic*),leftmargin=*]
\item \textbf{No independence result exists.} RH has never been shown
independent of PA, ZFC, or any standard theory. ``Open'' is not ``undecidable.''
A proof could arrive by a method outside every route catalogued in Diagnostic
thesis \ref{prin:mismatch} --- that catalogue is a description of \emph{past}
attempts, not a closed classification of all possible ones.
\item \textbf{If RH were independent of a sound theory, it would be true.} RH
has known $\Pi_1$ equivalents --- e.g.\ the Lagarias \cite{Lagarias} form of
Robin's inequality, an arithmetic statement whose failure at some $n$ would be
a finitely checkable counterexample. For any sound theory strong enough to
verify such counterexamples, independence of that $\Pi_1$ formulation would
imply it has no counterexample, i.e.\ would force RH \emph{true}. A true
$\Pi_1$ statement cannot be ``false-but-unprovable.'' Even the unprovability
door opens onto truth, not onto an unknowable void.
\item \textbf{The void is extensively known} (\S\ref{sec:character}): counted
to height $T$, at least $2/5$ on the line, $\sim10^{13}$ verified, statistics
classified. ``You can never know anything about it'' is simply false. What is
open is one uniform statement, not the object's knowability.
\end{enumerate}

\begin{tcolorbox}[colback=red!2,colframe=red!50!black,arc=2mm,
  title={\textbf{What this paper does and does not claim}},
  fonttitle=\bfseries,breakable]
\textbf{Claims (\TAGthm/\TAGver/\TAGev):} the zeros are an absence; the primes
are reconstructed from them; the void is timeless/observer-free and formless in
exact senses; and a principled \emph{diagnosis} of why construction stalls.\\[2pt]
\textbf{Does \emph{not} claim:} that RH is false, that RH is unprovable, that
RH ``can never be proven,'' or that the void is unknowable. Those are
\emph{not} consequences of anything above, and the mathematics does not
support them.
\end{tcolorbox}

\section{Scope discipline}\label{sec:discipline}

The mathematics here uses only L-functions, their zeros, the explicit formula,
and random-matrix classification. \textbf{No observer, no consciousness, no
``self,'' no purpose, and no deity enters any definition, and none is needed.}
The essay's defensible reach is the following ladder; the line marks exactly
where mathematics ends and a reader's interpretation begins.

\begin{center}
\begin{tcolorbox}[colback=gray!3,colframe=gray!55!black,arc=2mm,width=0.95\textwidth,
  title={\textbf{The ladder, and where it stops}},fonttitle=\bfseries]
\ttfamily\small
the void \hfill the zeros exist and are infinite \;[theorem]\\
generation \hfill primes $=$ the void's harmonics (explicit formula) \;[theorem]\\
formlessness \hfill GUE under restricted support; matched in tests \;[theorem(restr.)$+$ev.]\\
necessity \hfill determined once the integers are fixed \;[determinacy]\\
intelligibility \hfill in fact computed, to large height \;[fact]\\
eternity \hfill not in space/time; same for all; by proof, not faith \;[exact sense]\\[2pt]
\textmd{\rule{0.9\textwidth}{0.4pt}}\\[2pt]
$\blacktriangle$ above: each rung at the strength actually proved --- undeniable as written\\
$\blacktriangledown$ below: the reader's, and the text does not enter it\\[2pt]
life \;$\cdot$\; consciousness \;$\cdot$\; meaning \;$\cdot$\; the sacred \;$\cdot$\; God \;$\cdot$\; a presence
\end{tcolorbox}
\end{center}

\begin{remark}[What ``undeniable'' requires, \TAGphil]\label{rem:undeniable}
Every rung above the line is stated at \emph{exactly} the strength that has
been proved. Two are theorems (the void, generation) and one a plain fact
(intelligibility); one (formlessness) is a theorem only under restricted
support, plus numerical evidence (Prop.\ \ref{prop:formless}); and two
(necessity, eternity) hold in a precisely fixed sense --- the zeros are
determined once the integers are fixed --- rather than by a substantive
theorem. An earlier draft listed a ``self-modelling (Gödel)'' rung as a
theorem; it is not proved in this paper and has been removed. This
is the \emph{only} honest route to ``undeniable'': one reaches it by
\emph{trimming each claim to its proven core}, never by strengthening the
language to sound stronger. A rung inflated past its proof is the first thing a
referee strikes out --- and it takes the credible rungs down with it. The
ladder is undeniable because, and only because, nothing on it reaches past what
is shown.
\end{remark}

Consistent with the programme's discipline (``$X$ exists $\Leftrightarrow
\mathcal{C}(X)=X$''; no ``primes are alive''; no ``universe is conscious''),
the items below the line are not claimed, not implied, and not required. The
restraint is deliberate: a description that points only at what is true is the
one a reader can trust when deciding, for themselves, what it means.

\section*{Conclusion}
The Riemann zeros are best understood as a \emph{generative absence}: a void,
in the exact sense of a vanishing locus realised spectrally as missing lines,
from which --- by the explicit formula --- the prime structure is
reconstructed as interference. The void is timeless, uncreated, and the same
for everyone, and its statistics are those of the unitary universality class
(proven under restricted support, matched in every test). Our methods of proof
build structure with shape; the void is an absence with
no shape beyond the arithmetic that pins it; and that mismatch is, on the
honest diagnosis offered here, why every construction has stalled at the same
wall. But hardness is not impossibility: nothing here shows RH cannot be
proved, and a great deal is known about the void. What the mathematics gives
is a true and bounded thing --- an eternal, observer-free, generative absence,
fully real and not taken on faith. What anyone makes of it past that line is
theirs to decide.

\begin{thebibliography}{9}
\bibitem{Connes} Connes, A. \emph{Trace formula in noncommutative geometry
and the zeros of the Riemann zeta function.} Selecta Math.\ 5 (1999), 29--106.
\bibitem{Montgomery} Montgomery, H.L. \emph{The pair correlation of zeros of
the zeta function.} Proc.\ Sympos.\ Pure Math.\ 24 (1973), 181--193.
\bibitem{RudnickSarnak} Rudnick, Z., Sarnak, P. \emph{Zeros of principal
L-functions and random matrix theory.} Duke Math.\ J.\ 81 (1996), 269--322.
\bibitem{Odlyzko} Odlyzko, A.M. \emph{On the distribution of spacings between
zeros of the zeta function.} Math.\ Comp.\ 48 (1987), 273--308.
\bibitem{Conrey} Conrey, J.B. \emph{More than two fifths of the zeros of the
Riemann zeta function are on the critical line.} J.\ reine angew.\ Math.\ 399
(1989), 1--26.
\bibitem{Platt} Platt, D.J. \emph{Isolating some non-trivial zeros of zeta.}
Math.\ Comp.\ 86 (2017), 2449--2467.
\bibitem{Gourdon} Gourdon, X. \emph{The $10^{13}$ first zeros of the Riemann
zeta function, and zeros computation at very large height.} (2004), preprint.
\bibitem{Wigner} Wigner, E.P. \emph{Characteristic vectors of bordered
matrices with infinite dimensions.} Ann.\ of Math.\ 62 (1955), 548--564.
\bibitem{Stockmann} St\"ockmann, H.-J. \emph{Quantum Chaos: An Introduction.}
Cambridge University Press, 1999.
\bibitem{Lagarias} Lagarias, J.C. \emph{An elementary problem equivalent to
the Riemann hypothesis.} Amer.\ Math.\ Monthly 109 (2002), 534--543.
\bibitem{Dyson} Dyson, F.J. \emph{Birds and frogs.} Notices Amer.\ Math.\ Soc.\
56 (2009), 212--223.
\bibitem{ConnesConsani} Connes, A., Consani, C. \emph{Schemes over $\F_1$ and
zeta functions.} Compositio Math.\ 146 (2010), 1383--1415.
\bibitem{Deninger} Deninger, C. \emph{Some analogies between number theory and
dynamical systems on foliated spaces.} Proc.\ ICM Berlin 1998, Vol.\ I, 163--186.
\bibitem{IK} Iwaniec, H., Kowalski, E. \emph{Analytic Number Theory.} AMS
Colloq.\ Publ.\ 53, 2004.
\bibitem{BridgePaper} Smart, L. \emph{From Dirac Solutions to Physical
Reality: A Crystallisation-Based Selection Architecture.} VFD programme
(bridge paper); the Born-rule recovery is stated as conjectured, not proven.
\bibitem{ProbGeom} Smart, L. \emph{Probability as Constraint Geometry.} VFD
programme (Paper XXX).
\end{thebibliography}

\end{document}
